\chapter{Postup riešenia zadaných úloh}
\label{sec:Riesenie}

% TODO: Napísať detailný popis riešenia (15-25 strán)
% Podľa zadania: bod 3 - zvolený postup riešenia
% Toto je JADRO práce - najdlhšia kapitola

Táto kapitola detailne popisuje postup riešenia jednotlivých zadaných úloh.

% ============================================
\section{Úloha 1: Názov}
\label{sec:ukol1}

\subsection{Zadanie a ciele}
% Čo bolo cieľom úlohy?

\subsection{Analýza problému}
% Ako ste problém analyzovali?

\subsection{Návrh riešenia}
% Aký bol navrhnutý postup?

\subsection{Implementácia}
% Ako ste riešenie implementovali?
% Ukážky kódu, screenshoty, diagramy

% Príklad vloženia kódu:
% \begin{lstlisting}[label=src:example1, caption={Popis kódu}]
% // váš kód tu
% \end{lstlisting}

% Príklad vloženia obrázka:
% \begin{figure}[htbp]
%     \centering
%     \includegraphics[width=0.8\textwidth]{Figures/screenshot.png}
%     \caption{Popis obrázka}
%     \label{fig:example1}
% \end{figure}

\subsection{Testovanie}
% Ako ste riešenie testovali?

\subsection{Výsledky a zhodnotenie}
% Aký bol výsledok? Čo sa podarilo/nepodarilo?

% ============================================
\section{Úloha 2: Názov}
\label{sec:ukol2}

% Rovnaká štruktúra ako Úloha 1

\subsection{Zadanie a ciele}

\subsection{Analýza problému}

\subsection{Návrh riešenia}

\subsection{Implementácia}

\subsection{Testovanie}

\subsection{Výsledky a zhodnotenie}

% ============================================
% Pridať ďalšie sekcie pre ďalšie úlohy...

\endinput
